On breast cancer Wisconsin with LR and GB, none of reweighting, random oversampling or SMOTE gave a reliable gain in AUROC or AUPRC at any prevalence. Brier score was higher with a correction in 11 of 12 cells at 1:20 and 1:50 but without statistical support, and the absolute calibration intercept was not consistently larger. What the corrections did for LR was shift predicted risk upward: the signed intercept fell in every cell and went from underestimation to overestimation at 1:5--1:50. At the 0.5 cut-off this raised LR sensitivity at every prevalence and balanced accuracy at 1:5--1:50; threshold moving, which cannot change predicted risks, raised both more at 1:20 and 1:50. For GB neither helped reliably at 1:20 and 1:50, and the analytic prior-shift offset did not repair the corrected models. With one easy dataset and 10 seeds the advice is modest: report the signed calibration intercept with discrimination, and try the threshold before resampling when only the operating point matters.
